% ===================================================================
%  جدول نمبر 7 — جاڑے میں گاؤں۔ بہار میں کھیت کا گھر۔
%
%  Traduction, faite en 2026, en ourdou d'aujourd'hui, sur l'ido de
%  texto/io. Regles de la colonne : voir 10-jadval-01.tex. Deux
%  scenes, deux numerotations : les cles portent c1 et c2. Ecrit avec
%  outils/ordo.py sous les yeux.
%
%  LE MARRON CHAUD (42), ET L'OURDOU EST LA DEUXIEME LANGUE DU RELEVE
%  A L'AVOIR. C'est le trou le mieux documente de tout le travail :
%  le telougou, le marathe, l'egyptien et le tamoul ont tous du
%  periphraser, parce que le chataignier ne pousse pas chez eux ; le
%  coreen l'avait, avec le mot ET la scene. L'ourdou l'a aussi. Le
%  chataignier du Cachemire et de la frontiere donne la شاہ بلوط, et
%  le گرم چھلی — les marrons grilles vendus dans la rue en hiver —
%  est une chose de Lahore et de Peshawar. Cinq colonnes sans, deux
%  avec, et la coupure suit toujours la carte.
%
%  LE CORDONNIER (7) ET SA BLAGUE, ET L'OURDOU N'A PAS LE PROBLEME
%  DES AUTRES. Le coreen avait du refuser 갖바치, le tamoul son nom
%  de caste, le telougou son కసాయి : tous des mots justes devenus des
%  injures. L'ourdou dit موچی, qui est le nom du metier et un nom de
%  famille ordinaire au Pendjab — comme le चांभार marathe, que la
%  colonne marathe avait pu garder. Meme question, deux reponses
%  opposees, et c'est la langue vivante qui tranche a chaque fois,
%  pas la symetrie.
%
%  LA FERME EST ENTIERE EN OURDOU : ہل la charrue, نالی le sillon,
%  سہاگا la herse, کدال la pioche, چرخہ le rouet, کنواں le puits,
%  طویلہ l'etable, ریوڑ le troupeau, گوبر le fumier, کلغی la crete,
%  مور le paon, بچھڑا le veau, بچھیرا le poulain, چوزہ le poussin.
%  Et l'ourdou separe les sexes par un mot entier la ou l'ido met
%  -ulo et -ino : بیل le taureau contre گائے la vache, مرغا le coq
%  contre مرغی la poule.
%
%  L'OIE (69) NE MANQUE PAS ICI. Le tamoul avait du ecrire
%  பெருவாத்து, le grand canard, faute d'oiseau de basse-cour ;
%  l'ourdou a راج ہنس pour l'oie et بطخ pour le canard, deux mots
%  distincts et courants. Le Pendjab eleve les deux.
%
%  LE CYPRES (22) A CHANGE DE MOT APRES COUP, ET LE TABLEAU 9 EN EST
%  LA CAUSE. La colonne avait d'abord ecrit صنوبر, qui se dit bien
%  d'un conifere mais qui est en ourdou le sapin ; le cypres est
%  سرو, mot persan et mot de la poesie ourdoue pour l'arbre droit
%  des cimetieres. La faute n'est apparue qu'au tableau 9, ou l'ido
%  demande le SAPIN a deux endroits (11 et 31) : il fallait alors
%  deux mots, et le premier etait pris. Corrige ici. Une colonne se
%  relit d'un tableau a l'autre, pas seulement bloc a bloc.
%
%  LE BLOC 12 DE LA SECONDE SCENE FAIT TREBUCHER LA TROISIEME
%  COLONNE DE SUITE, ET LA PARADE N'EST PAS LA MEME. L'ido enchaine
%  « fromaji (64), qui gutifas de la tabulo pozita sur la kuvego
%  (63) » : une relative posee APRES son antecedent. Le coreen et le
%  tamoul n'ont pas de relative posee apres, seulement un participe
%  place avant, si bien que la cuve remontait devant les fromages et
%  que renvoji.py criait ; le tamoul avait du couper la phrase en
%  deux et reprendre « ces fromages-la ». L'ourdou n'a pas besoin de
%  couper : جو ... ہیں est une relative ordinaire de l'ourdou ecrit,
%  heritee du persan, et elle se pose apres son antecedent comme
%  celle de l'ido. Meme piege, meme diagnostic, et une sortie plus
%  courte parce que la langue a l'outil. A retenir pour les colonnes
%  qui restent : avant de couper une phrase en deux, verifier si la
%  langue possede une relative postposee.
%
%  TRENTE-DEUX PAIRES. Le tableau se partage en deux et le partage se
%  lit dans le compte : la premiere scene raconte des visites et
%  enumere des artisans, la seconde POSE la basse-cour objet par
%  objet. C'est la seconde qui coute, comme au coreen et au tamoul.
% ===================================================================

%%K t07-noto-1 noto
\VUnotes{143.2mm}{%
(*) کھانے کی تفصیل کے لیے جدول 6 اور 13 دیکھیے۔%
}
%%K t07-apar-1 apar
\VUsaut{4.0mm}
\VUcentre{13.2pt}{90}{دوسرا سلسلہ}
\VUsaut{3.0mm}
\VUfilet{16mm}
\VUsaut{5.6mm}
\VUtitre{15.0pt}{60}{جدول نمبر 7}
\VUsaut{3.0mm}

%%K t07-tit-1 sub
\VUcentre{12.0pt}{0}{\textit{پہلا منظر۔}}
\VUsaut{3.0mm}

%%K t07-tit-2 sub
\VUpk{10.2pt}{40}{جاڑے میں گاؤں۔}
\VUsaut{4.0mm}

%%K t07-c1-01-1 p
1. --- نئے سال کی چھٹیوں میں، میں اپنے والد کے ساتھ نو اُوربو نامی
ایک چھوٹے گاؤں گیا ؛ وہاں اُنھیں کچھ کاروباری کام تھے۔

%%K t07-c1-02-1 p
2. --- ہم نے شہر اور اُس قصبے کے بیچ چلنے والی \VUgras{ڈاک گاڑی}
\textsuperscript{(11)} لی ؛ مگر بہت سردی تھی، اِس لیے اُس کم آرام دہ
گاڑی کا سفر زیادہ خوشگوار نہ رہا۔

%%K t07-c1-03-1 p
3. --- اسی لیے جب ہم نے \VUgras{گاڑی بان} کو \textit{سفید ریچھ} نامی
\VUgras{سرائے} \textsuperscript{(2)} کے سامنے چوک میں اپنی گاڑی
روکتے دیکھا تو ہمیں سچی خوشی ہوئی۔ \VUgras{سرائے والا}
\textsuperscript{(4)} اپنے گھر کی دہلیز پر کھڑے اطمینان سے اپنا
\VUgras{حقہ} پیتے ہمارا انتظار کر رہے تھے۔

%%K t07-c1-03-2 p
اُنھوں نے ہمیں اندر بلایا ؛ چولھے میں چٹختی ٹہنیوں کی آگ کے سامنے
بیٹھنے کے لیے میں نے دوسری بار کہلوانے کا انتظار نہ کیا۔ تب میں نے
اُس تیز \VUgras{شمالی ہوا} کا خوب مذاق اڑایا جو پرانے
\VUgras{سائن بورڈ} \textsuperscript{(3)} کو چرچرا رہی تھی اور چمنی سے
نکلتے \VUgras{دھوئیں} \textsuperscript{(1)} کو کالے بھنوروں میں اڑا
لے جا رہی تھی۔

%%K t07-c1-04-1 p
4. --- ناشتے کا وقت تھا۔ مزید انتظار کیے بغیر ہم نے وہ سادہ مگر مزے
دار کھانا خوب کھایا جو ہمیں پیش کیا گیا (*)۔

%%K t07-tit-3 sub
\VUpk{10.2pt}{40}{کچھ ملاقاتیں۔}
\VUsaut{3.0mm}

%%K t07-c1-05-1 p
5. --- ناشتے کے بعد میں اپنے والد کے ساتھ گیا، جو بیمہ کے ایجنٹ ہیں،
اُن کے گاہکوں کے پاس۔ پہلے \VUgras{لوہار} \textsuperscript{(5)} کے
پاس گئے ؛ وہ سرائے کے قریب رہتے ہیں۔ وہ ایک گھوڑے کو نعل لگانے میں
مصروف تھے۔ میں نے اُن کے مددگار کی طاقت کو سراہا ؛ وہ اپنے چمڑے کے
تہبند پر گھوڑے کا \VUgras{سُم} مضبوطی سے تھامے ہوئے تھا ؛ اُسی وقت
لوہار زور دار ہتھوڑے کی ضربوں سے \VUgras{نعل} \textsuperscript{(6)}
جما رہے تھے۔

%%K t07-c1-06-1 p
6. --- پھر ہم گاؤں کے \VUgras{موچی} \textsuperscript{(7)} کے پاس
گئے، بوڑھے یاکوبس کے پاس۔ اگرچہ اُن کا سائن بورڈ ایک بہت خوبصورت
\VUgras{بوٹ} ہے، وہ ایک پرانے سوراخ دار جوتوں کے جوڑے میں نیا تلا
لگانے پر ہی مطمئن تھے۔

%%K t07-c1-06-2 p
اُس بوڑھے نے ہنستے ہوئے ہم سے کہا : « شہر میں میں « جوتے بنانے والا »
تھا اور میرے کوئی گاہک نہ تھے۔ یہاں میں « جوتے مرمت کرنے والا » ہوں
اور کام بہت ہے۔ »

%%K t07-c1-07-1 p
7. --- اُنھوں نے ہمیں اپنے پڑوسی \VUgras{حجام} \textsuperscript{(8)}
کے بارے میں بتایا ؛ اُس کے گاہک خوش نہیں۔ اُس کے \VUgras{استرے}
ہمیشہ کند رہتے ہیں اور اُس کی \VUgras{قینچیاں} کبھی اچھی طرح نہیں
کاٹتیں۔

مگر وہی گاؤں کا واحد حجام ہے، اِس لیے مجبوراً اُسی سے حجامت اور بال
کٹوانے پڑتے ہیں۔ اور وہ کنجوس اور بدمزاج آدمی بھی ہے۔

%%K t07-c1-08-1 p
8. --- خوش قسمتی سے باقی \VUgras{پڑوسی} اُس جیسے نہیں۔ مثلاً
\VUgras{گاڑی بنانے والا} \textsuperscript{(9)} نیک آدمی ہے، محنتی اور
مددگار۔ اُس سے گاڑی مرمت کرانا خوشی کی بات ہے۔ اُس کا کام ہمیشہ
مکمل ہوتا ہے۔

%%K t07-c1-09-1 p
9. --- بوڑھے یاکوبس نے \VUgras{زین ساز} \textsuperscript{(10)} کی
شکایت بھی نہیں کی ؛ کام کے جلدی میں ہونے پر وہ کبھی کبھی اُس کے
\VUgras{ساز} سینے میں مدد کرتے ہیں۔

%%K t07-c1-09-2 p
میرے والد نے اُس سے اُس کے خاندان کے بارے میں پوچھا : « سب خیریت سے
ہیں۔ میری پوتی \VUgras{اسکول} \textsuperscript{(12)} میں ہے۔ اُس کی
\VUgras{استانی} \textsuperscript{(13)} اُس سے بہت خوش ہیں ؛ وہ اچھی
\VUgras{طالبہ} \textsuperscript{(14)} ہے۔ وہ میرے نالائق بیٹے جیسی
نہیں ؛ وہ صرف کھیلنے کا سوچتا ہے۔ \VUgras{استاد}
\textsuperscript{(15)} اُسے کچھ بھی نہیں سکھا پاتے۔ »

%%K t07-tit-4 sub
\VUsaut{3.0mm}
\VUpk{10.2pt}{40}{گاؤں کی باتیں۔}
\VUsaut{3.0mm}

%%K t07-c1-10-1 p
10. --- پھر اُس بوڑھے نے ہمیں گاؤں کی خبریں سنائیں۔ نیا اسکول بننے
والا ہے، کیونکہ \VUgras{بلدیہ کے گھر} میں بنا ہوا اسکول بہت چھوٹا پڑ
گیا۔ یہ آسان ہے، کیونکہ \VUgras{چکی والے} کے باغ کا ایک حصہ خرید لینا
کافی ہو گا۔ اُس بے چارے آدمی کے لیے یہ بہت فائدہ مند ہو گا۔ سردی کی
وجہ سے \VUgras{چکی} \textsuperscript{(16)} کے \VUgras{پاٹ} اب گندم
نہیں پیس سکتے ؛ کیونکہ \VUgras{پہیے} \textsuperscript{(17)} کو
\VUgras{برف} \textsuperscript{(25)} نے جما کر روک دیا ہے ؛ وہ برف
\VUgras{ندی} \textsuperscript{(23)} کی ہے۔

%%K t07-c1-11-1 p
11. --- « عمارتوں کی بات کریں تو، ہمارا نیا \VUgras{کلیسا}
\textsuperscript{(18)} آپ نے دیکھا؟ برف کی چادر کے نیچے وہ بہت
خوبصورت لگتا ہے۔ \VUgras{کوے} \textsuperscript{(19)} اپنے پرانے گھنٹہ
گھر کو اب نہیں پہچانتے، اور \VUgras{قبرستان} \textsuperscript{(20)}
کے بڑے درخت پر اداس بیٹھے ہیں۔ »

%%K t07-c1-12-1 p
12. --- قبرستان کے اُس خیال نے موچی کو وہ لوگ یاد دلائے جنھیں میرے
والد جانتے تھے اور جو پچھلے سال مر گئے۔ « کتنی \VUgras{قبریں}
\textsuperscript{(21)}، میرے اچھے صاحب، \VUgras{سرو}
\textsuperscript{(22)} کے نیچے باقیوں میں جا ملیں! موت نے ہمارے پرانے
نوٹری کو کاٹ ڈالا۔ سب گاؤں والے اُن کے \VUgras{جنازے} میں شریک
ہوئے۔ »

%%K t07-c1-13-1 p
13. --- « یہ رہے اُن کے جانشین ؛ ندی پر پھسل رہے ہیں۔ ابھی آئے تھے
کہ میں اُن کے ٹوٹے ہوئے \VUgras{برف کے جوتے} \textsuperscript{(26)}
کا تسمہ مرمت کر دوں۔ »

%%K t07-c1-14-1 p
14. --- « ندی کے کنارے یہ رہا نیا \VUgras{دیہاتی چوکیدار}
\textsuperscript{(27)}۔ وہ سابق سپاہی ہے ؛ گاؤں کے شرارتی لڑکوں کو
ڈراتا ہے۔ اُس کے \VUgras{پنڈلی پوش} \textsuperscript{(29)} اور اُس کی
\VUgras{وردی کی ٹوپی} \textsuperscript{(28)} دیکھتے ہی وہ بھاگ جاتے
ہیں۔ »

%%K t07-c1-15-1 p
15. --- « آہا! یہ رہے بوڑھے یوسف ؛ نان بائی کے لیے \VUgras{لکڑی کے
گٹھوں کی چھوٹی گاڑی} \textsuperscript{(30)} ہانک رہے ہیں۔ --- سچ ہے،
کہا میرے والد نے ؛ اُن کے \VUgras{خچروں} \textsuperscript{(31)} کی
جوڑی میں پہچانتا ہوں۔ مجھے اُن سے بات کرنی ہے ؛ اِس موقع سے فائدہ
اٹھاتا ہوں۔ پھر ملیں گے، یاکوبس بابا۔ »

%%K t07-c1-16-1 p
16. --- جس وقت ہم باہر نکلے، گاؤں کے لڑکے اسکول سے نکل کر دوڑتے ہوئے
\VUgras{پھسلنے کی جگہ} \textsuperscript{(33)} کی طرف لپکے — \VUgras{گر}
\textsuperscript{(34)} کر ہاتھ پاؤں توڑ لینے کے خطرے کی پروا کیے
بغیر۔ اُن میں سے ایک نے گاڑی بنانے والے سے اپنی \VUgras{پھسل گاڑی}
\textsuperscript{(32)} لی، اُس میں اپنے ایک ساتھی کو بٹھایا، اور
دوڑتا ہوا چل دیا۔

%%K t07-c1-17-1 p
17. --- کچھ اور \VUgras{برف کے ڈھیر} \textsuperscript{(37)} کی طرف
لپکے ؛ وہ \VUgras{پُل} \textsuperscript{(40)} کے قریب ہے۔ کل جو
\VUgras{برف کا مجسمہ} \textsuperscript{(39)} اُنھوں نے بنایا تھا، اُس
پر گولہ باری شروع کر دی ؛ اُس مجسمے کو اُنھوں نے ٹوپی، حقہ اور کندھے
پر لٹکتا بستہ پہنایا تھا۔

%%K t07-c1-18-1 p
18. --- جما دینے والی ہوا اور سردی کی وہ زیادہ پروا نہیں کرتے۔ اُن
میں سے اکثر کے پاس \VUgras{گلوبند} \textsuperscript{(35)} تک نہیں۔
\VUgras{برف کے گولوں} \textsuperscript{(38)} کی لڑائی کے بعد پھر گرم
ہونے کے لیے اُنھیں مٹھی بھر بہت گرم \VUgras{شاہ بلوط}
\textsuperscript{(42)} کافی ہیں ؛ وہ اُنھیں بوڑھی \VUgras{بیچنے والی}
یاکوبینا سے خریدتے ہیں ؛ وہ اپنی بہت باریک \VUgras{پھٹی چادر}
\textsuperscript{(43)} کے نیچے سردی سے کانپ رہی ہے۔

%%K t07-tit-5 sub
\VUsaut{3.0mm}
\VUcentre{12.0pt}{0}{\textit{دوسرا منظر۔}}
\VUsaut{3.0mm}

%%K t07-tit-6 sub
\VUpk{10.2pt}{40}{بہار میں کھیت کا گھر۔}
\VUsaut{4.0mm}

%%K t07-c2-01-1 p
1. --- جاڑے کی سب خوشیوں کے باوجود ہم \VUgras{بہار} کی واپسی کا گہری
خوشی سے استقبال کرتے ہیں۔

%%K t07-c2-01-2 p
مئی کے مہینے میں، شام کے وقت، کھیت کا صحن کیسا خوشگوار منظر لگتا ہے!

%%K t07-c2-02-1 p
2. --- افق پر \VUgras{سورج} \textsuperscript{(1)} آہستہ آہستہ ڈھل رہا
ہے ؛ \VUgras{دیہات} \textsuperscript{(3)} کو اپنی ارغوانی
\VUgras{کرنوں} \textsuperscript{(2)} سے بھر رہا ہے۔ محنتی
\VUgras{ہل چلانے والا} \textsuperscript{(6)} اپنا \VUgras{کھیت}
\textsuperscript{(4)} جوتنے میں جلدی کر رہا ہے۔

%%K t07-c2-02-2 p
اپنے \VUgras{ہل} \textsuperscript{(5)} سے — وہ ایک زور دار
\VUgras{گھوڑے} \textsuperscript{(7)} سے جتا ہے — وہ \VUgras{نالی}
\textsuperscript{(8)} کھینچتا ہے ؛ اُس میں \VUgras{بیج بونے والا}
\textsuperscript{(9)} مٹھی بھر بھر کر \VUgras{بیج} ڈالے گا ؛ اُس سے
سنہری بالیاں نکلیں گی۔

%%K t07-c2-03-1 p
3. --- \VUgras{گھاس کاٹنے والوں} \textsuperscript{(11)} نے اپنی
درانتیوں سے چراگاہ کی گھاس کاٹی ؛ سکھانے والوں نے \VUgras{سوکھی
گھاس} \textsuperscript{(10)} کو \VUgras{کانٹوں}
\textsuperscript{(12)} اور پنجوں سے الٹ پلٹ کر سکھایا ؛ اور یہ رہے
وہ واپس آ رہے ہیں۔

%%K t07-c2-03-2 p
کچھ بیلوں کی کھینچی بھاری گاڑی کے آگے چل رہے ہیں ؛ اپنا اوزار کندھے
پر اٹھائے ہوئے ہیں۔ کچھ اور، زیادہ سست، مہکتی سوکھی گھاس پر آرام سے
بیٹھ گئے ہیں۔

%%K t07-c2-04-1 p
4. --- گزرتے ہوئے وہ \VUgras{مالی} \textsuperscript{(16)} کو دوستانہ
سلام کرتے ہیں ؛ وہ \VUgras{پھلوں کے باغ} \textsuperscript{(13)} میں
\VUgras{پھل دار درختوں} کی کانٹ چھانٹ اور \VUgras{ناشپاتی کے درختوں}
\textsuperscript{(14)} کی پیوند کاری میں لگا ہے۔ \VUgras{آلوچے کے
درخت} \textsuperscript{(15)} پھولنے لگے ہیں ؛ اور خوب گوبر ڈالے
\VUgras{سبزیوں کے باغ} \textsuperscript{(17)} میں سبزیاں، بند گوبھیاں
اور سلاد پہلے ہی اچھی حالت میں ہیں۔

%%K t07-c2-04-2 p
چھوٹی دیوار پر خوبصورت \VUgras{مور} \textsuperscript{(18)} اپنی دم
پھیلا رہا ہے ؛ اپنے بہت خوبصورت پروں کو سراہوانے کے لیے۔

%%K t07-tit-7 sub
\VUpk{10.2pt}{40}{کھیت کا صحن۔}
\VUsaut{3.0mm}

%%K t07-c2-05-1 p
5. --- \VUgras{اصطبل} \textsuperscript{(19)} کے دروازے کھلے ہیں ؛
چارے کی کھرلی اور \VUgras{بچھونے کی گھاس} \textsuperscript{(23)} نظر
آتی ہیں ؛ وہ \VUgras{گھوڑی} \textsuperscript{(21)} کی ہیں ؛ اُس گھوڑی
کو \VUgras{سائیس} \textsuperscript{(20)} نے ابھی کھولا ہے۔ لمبی
ٹانگوں سے بہت مزاحیہ لگتا \VUgras{بچھیرا} \textsuperscript{(22)}
کلیلیں کرتا اپنی ماں کے پیچھے چل رہا ہے۔

%%K t07-c2-06-1 p
6. --- \VUgras{کنویں} \textsuperscript{(29)} کے قریب \VUgras{گائیں}
\textsuperscript{(25)} لکڑی کے \VUgras{کھرلے} \textsuperscript{(26)} میں
پانی پینے جا رہی ہیں ؛ وہ کھرلا اُن کے لیے پانی پینے کی جگہ ہے۔
\VUgras{بچھڑا} \textsuperscript{(24)} اِس موقع سے فائدہ اٹھا کر دودھ
پی رہا ہے ؛ \VUgras{بیل} \textsuperscript{(27)} چارے کے گودام کی اچھی
خوشبو سونگھ کر خوشی سے ڈکارتا ہے، اور زور دار \VUgras{سینگوں}
\textsuperscript{(28)} والا اپنا سر فخر سے اٹھاتا ہے۔

%%K t07-c2-07-1 p
7. --- سب کام کر رہے ہیں۔ \VUgras{خادمہ} \textsuperscript{(32)}
کنویں کی \VUgras{رسّی} \textsuperscript{(31)} ہاتھ سے تھامے ہوئے ہے۔
مویشیوں کو پلانے کے لیے \VUgras{ڈول} \textsuperscript{(30)} سے پانی
کھینچ رہی ہے۔ \VUgras{بھوسے کے گودام} \textsuperscript{(33)} کے قریب
ایک آدمی بھوسا سمیٹ رہا ہے ؛ اور \VUgras{گھر} \textsuperscript{(34)}
کے دروازے کے سامنے بوڑھی \VUgras{سوت کاتنے والی}
\textsuperscript{(35)} اپنے چرخے اور اپنی \VUgras{تکلی} سے پرانے
زمانے کی طرح \VUgras{السی} یا \VUgras{سن} کات رہی ہے۔

%%K t07-tit-8 sub
\VUsaut{3.0mm}
\VUpk{10.2pt}{40}{کاشتکار۔}
\VUsaut{3.0mm}

%%K t07-c2-08-1 p
8. --- \VUgras{کاشتکار} \textsuperscript{(36)} یہ سارے کام چلاتا ہے
اور اپنے ملازموں کو ہدایتیں دیتا ہے۔ وہ \VUgras{سہاگے}
\textsuperscript{(40)} اور \VUgras{کدال} \textsuperscript{(39)} کو چھت
کے نیچے رکھنے کو کہتا ہے ؛ اُنھیں \VUgras{جھونپڑی}
\textsuperscript{(38)} کے پاس بھول گئے ہیں ؛ وہ جھونپڑی \VUgras{کتے}
\textsuperscript{(37)} کی ہے۔

%%K t07-c2-09-1 p
9. --- اُس کی آوازیں ایک \VUgras{کبوتر} \textsuperscript{(41)} کو ڈرا
دیتی ہیں ؛ وہ پورے پروں سے \VUgras{کبوتر خانے} \textsuperscript{(42)}
میں اڑ جاتا ہے۔ \VUgras{مرغیاں} \textsuperscript{(43)} بھی ڈر جاتی
ہیں ؛ وہ جلدی سے \VUgras{مرغی خانے} \textsuperscript{(44)} میں واپس
چلی جاتی ہیں۔

%%K t07-c2-10-1 p
10. --- \VUgras{بھیڑ خانے} \textsuperscript{(48)} کے سامنے، یہ رہا
بوڑھا \VUgras{چرواہا} \textsuperscript{(45)} ؛ وہ اپنا \VUgras{ریوڑ}
\textsuperscript{(49)} واپس لا رہا ہے۔ اپنی \VUgras{لاٹھی}
\textsuperscript{(46)} تھامے ہوئے — جو کبھی اُس سے جدا نہیں ہوتی،
اُس کی رنگ اڑی \VUgras{قمیص} \textsuperscript{(47)} کی طرح —، وہ
\VUgras{مینڈھوں} \textsuperscript{(50)}، \VUgras{بھیڑوں}
\textsuperscript{(53)}، \VUgras{میمنوں} \textsuperscript{(54)}،
\VUgras{بکریوں} \textsuperscript{(51)} اور \VUgras{بکری کے بچوں}
\textsuperscript{(52)} کو باڑے کی طرف لے جا رہا ہے۔ اُس کا وفادار
\VUgras{کتا} \textsuperscript{(55)} آرگس سب سے پیچھے آتا ہے ؛ اپنے
بھونکنے سے پیچھے رہ جانے والوں کو تیز کرتا ہے۔

%%K t07-c2-11-1 p
11. --- یہ سارا شور اور ہلچل اچھی \VUgras{دودھ دینے والی گائے}
\textsuperscript{(56)} کے سکون میں خلل نہیں ڈالتی ؛ وہ اپنی
\VUgras{گھنٹی} \textsuperscript{(57)} بجاتی ہے ؛ اُسی وقت
\VUgras{طویلے} \textsuperscript{(58)} کے دروازے کے سامنے اُسے دوہا جا
رہا ہے۔

%%K t07-c2-12-1 p
12. ایسا لگتا ہے کہ وہ سمجھتی ہے کہ اُسے دودھ دینا ہے ؛ تبھی تو
\VUgras{کاشتکار کی بیوی} \textsuperscript{(60)} \VUgras{بلوونے}
\textsuperscript{(61)} میں \VUgras{مکھن} بنا سکے گی، یا وہ بہترین
\VUgras{پنیر} \textsuperscript{(64)} جو \VUgras{بڑے ٹب}
\textsuperscript{(63)} پر رکھے تختے سے ٹپکتے ہیں۔

%%K t07-c2-13-1 p
13. --- سارا \VUgras{دودھ خانہ} \textsuperscript{(59)}
\VUgras{کھیت کا ملازم} \textsuperscript{(65)} بہت صاف رکھتا ہے ؛ اُس
نے ابھی \VUgras{دودھ کی ہانڈیاں} \textsuperscript{(62)} اُس بڑی
بالٹی میں دھوئی ہیں جو وہ ہاتھ میں اٹھائے ہوئے ہے۔

%%K t07-tit-9 sub
\VUsaut{3.0mm}
\VUpk{10.2pt}{40}{مرغی خانہ۔}
\VUsaut{3.0mm}

%%K t07-c2-14-1 p
14. --- اپنے موٹے لکڑی کے کھڑاووں کے ساتھ وہ کچھ بھاری چلتا ہے ؛
اِس لیے \VUgras{ٹرکی} \textsuperscript{(68)}، \VUgras{بطخوں}
\textsuperscript{(66)}، \VUgras{بطخوں کے نروں} \textsuperscript{(67)}
اور \VUgras{راج ہنسوں} \textsuperscript{(69)} کو بھگا دیتا ہے ؛ وہ
\VUgras{اپنی چونچ} \textsuperscript{(70)} خوب کھول کر بے چینی سے
چلاتے ہیں۔

%%K t07-c2-14-2 p
\VUgras{مرغا} \textsuperscript{(71)}، جو زیادہ لڑاکا ہے، اپنی سرخ
\VUgras{کلغی} \textsuperscript{(72)} اٹھاتا ہے، اپنی \VUgras{دم}
\textsuperscript{(73)} کے پر جھاڑتا ہے، اور گونجتی « کُکڑوں کوں »
سناتا ہے۔

%%K t07-c2-15-1 p
15. --- \VUgras{ماں مرغی} \textsuperscript{(74)} \VUgras{گوبر}
\textsuperscript{(81)} پر اپنے \VUgras{چوزوں} \textsuperscript{(75)}
کے ساتھ کرید رہی ہے۔ \VUgras{سؤر کا بچہ} \textsuperscript{(78)} اپنی
ماں کی طرف بھاگتا ہے — اُس بڑی \VUgras{سؤرنی}
\textsuperscript{(77)} کی طرف جو ہمیں سؤر خانے کے پاس نظر آتی ہے۔

%%K t07-c2-16-1 p
16. --- اپنے خانوں میں \VUgras{خرگوش} \textsuperscript{(79)} نرم
\VUgras{گھاس} \textsuperscript{(80)} شوق سے کھا رہے ہیں ؛ وہ گھاس
کاشتکار کی بیٹی اُن کے لیے لاتی ہے۔ یوآنینا اپنے خرگوشوں سے بہت
محبت کرتی ہے ؛ اور روز مرغی خانے میں تازہ \VUgras{انڈے}
\textsuperscript{(76)} جمع کرنے کے بعد اپنے ننھے دوستوں سے ملنے آتی
ہے۔

\VUsaut{6.0mm}
\VUfilet{22mm}
